\documentclass[
    style=paintings,
    paper=screen,
    blackslide,
    nopagebreaks,
    fleqn
]{powerdot}

\usepackage[brazil]{babel}
\usepackage[utf8]{inputenc}
\usepackage {graphicx}


\pdsetup{
lf=Ep1,
rf=Laboratório de Visão Computacional,
palette=GoldenGate,
theslide=slide~\arabic{slide},
list={itemsep=6pt},
%randomdots,
%dprop={dotstyle=ocircle,linewidth=.25pt},
%dmindots=10,dmaxdots=30,
%dbright=10,
}



\title{Busca em Imagens}
\author{Rafael Lopes \and Rodrigo Okada \and William Mizuta}
\date{
September, 2010
\vfill
\includegraphics[height=.1\slidewidth ]{img/logo_usp_b.ps}
\includegraphics[height=.1\slidewidth ]{img/logo_ime_b.ps}
}
\begin{document}
\maketitle

\begin{slide}{Outline}
\tableofcontents{}
\end{slide}

\section[slide=false]{Secao 1}
\begin{slide}{Introduction}
\begin{enumerate}[type=1]
\item The first thing;\pause
\item The second thing;\pause
\item The third thing;\pause
\item The final thing!
\end{enumerate}
\end{slide}


\section[slide=false]{Secao 2}

\begin{slide}{Objectives}
Teste de coisas com pausas:
\begin{itemize}[]
\item <2> The first;
\item <3> The second;
\begin{itemize}[type=1]
\item <4> One,
\item <5> Two,
\item <6> Three,
\end{itemize}
\item <7> The last one.
\end{itemize}
\end{slide}


\begin{slide}{Results}
  \begin{equation}\label{binomium}
      (a+b)^n=\sum_{k=0}^n{n\choose k}a^{n-k}b^k
        \end{equation}\pause
            Click the title of this slide to switch to the black slide.
\end{slide}


\section[slide=false]{Secao sem slide}
\begin{slide}{link back}
          We will prove formula (\ref{binomium}) on the blackboard.\\

\end{slide}



\begin{slide}{References}
\bibliographystyle{plain}
\bibliography{references}
\end{slide}


\end{document}
